\documentclass[11pt]{article}

\usepackage[utf8]{inputenc}
\usepackage[T1]{fontenc}
\usepackage{lmodern}
\usepackage{geometry}
\geometry{margin=1in}
\usepackage{setspace}
\usepackage{amsmath,amssymb}
\usepackage{booktabs}
\usepackage{graphicx}
\usepackage{hyperref}
\usepackage{natbib}
\usepackage{enumitem}

\title{Failure-by-Language: Reproducible Evaluation of Python\texttt{-}like Confusion Language in LLM Code Generation}
\author{Jaeyeong CHOI\\DGIST EECS\\\texttt{jaeyeong2022@dgist.ac.kr}}
\date{\today}

\begin{document}
\maketitle

\begin{abstract}
Recent studies on AI-assisted coding suggest a gap between short-term productivity and durable comprehension.
This paper proposes a reproducible benchmark for syntax-constrained code generation: LLMs are asked to write Python-like code under transformed-token rules.

We present an OCaml/CLI pipeline with lineage-aware schemas, run summary validation, and metric snapshot automation.
Our current empirical status is constrained by live API authentication issues, but fixture-based replay and failure-taxonomy analysis provide a controlled baseline for further model comparisons.
\end{abstract}

\section{Introduction}
Recent studies on AI-assisted coding consistently show a tension between short-term coding speed and durable understanding, especially for novices.\footnote{See structured review notes in \texttt{docs/research/literature/v1-curated.md}.}

This work studies a simple but strict question: when a model is instructed to write code in a Python-like language with transformed syntax, does it follow the new language rules, or revert to canonical Python conventions?

We define a reproducible pipeline in three layers:
\begin{enumerate}
  \item \textbf{Specification layer}: one-to-one token aliases and collision constraints.
  \item \textbf{Execution layer}: OCaml-based transform/validate/roundtrip tooling.
  \item \textbf{Measurement layer}: summary validation, metric generation, and schema checks.
\end{enumerate}

Our objective is not model ranking. It is to isolate where and how models fail under constrained language perturbation.

\section{Related Work}

\subsection{Code-Generation Benchmarks and LLM Priors}
HumanEval~\cite{chen2021codex} and MBPP~\cite{austin2021program} established that LLMs trained on large code corpora develop strong token-level priors: Python keywords appear in statistically predictable positions. More recent code LLMs (DeepSeek-Coder~\cite{guo2024deepseek}, WizardCoder~\cite{luo2023wizardcoder}) amplify these priors through specialized pretraining, while encoder-decoder models (CodeBERT~\cite{feng2020codebert}, CodeT5~\cite{wang2021codet5}) reinforce them via masked language modeling objectives.

\subsection{Instruction-Following in Code Generation}
Recent benchmarks---CodeIF~\cite{codeif2025} and MultiCodeIF~\cite{multicode2025}---reveal that LLMs underperform on code tasks that require strict adherence to developer-specified constraints, even when those constraints are explicitly stated. These works show a gap between general programming competence and instruction compliance, motivating our study of a more extreme case: instructions that directly conflict with Python NTP priors.

\subsection{Knowledge Conflicts: Parametric vs.\ Contextual}
A growing body of work studies how LLMs resolve conflicts between pretrained parametric knowledge and information given in context.
Chen et al.~\cite{chen2024instructhierarchy} propose an instruction hierarchy to prioritize privileged instructions; however, they focus on natural-language directives rather than programming-language semantic rules.
Xie et al.~\cite{xie2024conflict} and Wu et al.~\cite{wu2024contextparametric} show that LLMs often default to parametric knowledge when context contradicts it---especially under low surface-form overlap---which directly motivates our L3/L4 designs where syntax is identical to Python but semantics are inverted.
Wang et al.~\cite{wang2024ntp} characterize how NTP training shapes model behavior; their ``law of next-token prediction'' provides theoretical grounding for our observation that prior dominance scales with token predictability.

\subsection{Adversarial Perturbations in Code Generation}
PromptBench~\cite{zhu2023promptrobust} demonstrates LLMs are fragile under adversarial text perturbations at word and sentence levels.
Pan et al.~\cite{pan2024adversarial} show that semantically aligned adversarial code prompts cause up to 42\% accuracy drops.
Khan et al.~\cite{khan2024cheatingllm} specifically target programming assignment assessments, designing adversarial perturbations that reduce LLM correctness by 77\% to prevent academic dishonesty.
Our approach is complementary: rather than perturbation-based attack, we design a principled \emph{language family} in which LLM priors are systematically misaligned.

\subsection{LLM Impact on Programming Education}
Multiple empirical studies show LLMs reduce independent coding skill development~\cite{dakhel2023github,prather2023robots,lau2023learners}. Novak and McDanel~\cite{novak2025llmhw} explicitly study LLM-resistant assignment design, proposing contextual and personalized tasks; however, they do not characterize \emph{which linguistic features} drive LLM failure.
We address this gap by providing a four-level taxonomy grounded in NTP prior theory.

\subsection{Few-Shot Learning and Pattern Extraction}
In-context learning (ICL) allows LLMs to adapt to new tasks from examples~\cite{dong2022survey}.
Min et al.~\cite{min2022rethinking} show that labels in demonstrations matter less than expected---models rely on surface patterns rather than semantic reasoning.
Our L4 finding (pattern blindness) extends this: LLMs fail to extract \emph{semantic inversion} from worked examples, consistent with the view that ICL is primarily distributional rather than rule-inductive.

\subsection{Research Gap}
Prior adversarial work degrades generation quality without a principled design taxonomy. Prior education work proposes task-level heuristics without mechanism analysis. Prior knowledge-conflict work focuses on factual knowledge rather than programming-language semantics. Our contribution unifies these threads: a theory-grounded taxonomy (L1--L4) that characterizes LLM failure by \emph{type} (prior dominance vs.\ pattern blindness) and \emph{mechanism} (NTP prior conflict), enabling principled LLM-resistant language design.

\subsection{Instruction Alignment}
Alignment methods including RLHF~\cite{ouyang2022instruct} improve instruction following in natural-language tasks, but our results suggest semantic-level programming conflicts remain resistant even in aligned models (gpt-4o succeeds at L3 but fails at L4), pointing to a gap between alignment training and deep semantic inversion under implicit delivery.

\section{Method}
\subsection{Confusion Language Setup}
We map selected Python reserved tokens (e.g., \texttt{in}, \texttt{for}, \texttt{if}, \texttt{elif}, \texttt{return}, \texttt{def}) to alias phrases. To prevent decoding ambiguity, each alias must be unique (1:1 mapping).

\subsection{Validation and Transformation Pipeline}
The OCaml-based toolchain validates alias collisions, transforms source code into alias form, and checks roundtrip restoration. Supporting scripts generate summary and metric payloads with schema validation gates.

\subsection{Metrics}
We track:
\begin{itemize}[leftmargin=1.2em]
  \item ACR: Alias Compliance Rate
  \item PRR: Python Reversion Rate
  \item ESR: Execution Success Rate
  \item MFB: Manual Fix Burden
  \item LGP: Learning Gain Proxy
\end{itemize}

\section{Experiments}
\subsection{Run protocol}
Each run follows these steps:
\begin{enumerate}
  \item Prepare task-set metadata and prompt condition.
  \item Execute model batch or fixture replay.
  \item Validate summary payloads and metric snapshots, then archive in \texttt{docs/research/results/}.
\end{enumerate}

\subsection{Execution coverage}
Current coverage includes:
\begin{itemize}
  \item fixture summary replay with full artifact generation (JSON/CSV/Markdown) and schema validation;
  \item planner/state-code regression expansion;
  \item robust parsing output via compact JSON/JSONL/TSV/markdown formats.
\end{itemize}

\subsection{Failure conditions observed}
Live LLM batch calls are blocked by authentication failures (HTTP 401) in the current environment. Earlier local execution failures were additionally caused by missing environment dependencies (e.g., \texttt{dune}) and were resolved before this phase.

\subsection{Reproducibility protocol}
Every run is committed with provenance fields and lineage IDs (task-set, alias-set, prompt config, and run metadata) so that mismatches are auditable and not silently conflated with tool drift.

\section{Results}
Using the first fixture batch, we obtained the following initial metrics:
\begin{itemize}[leftmargin=1.2em]
  \item ACR = 0.333
  \item PRR = 0.667
  \item ESR = 0.667
  \item MFB = 65.0
  \item LGP = 0.667
\end{itemize}

These values indicate low alias compliance and relatively high reversion to Python priors. While this is only an initial fixture run, the direction is consistent with the central hypothesis of grammar-prior dominance.

\section{Discussion}
The current stage contributes a robust benchmark scaffold \emph{and} a meaningful real-call accumulation (480 samples), but still no passing cases under strict compliance.

\subsection{What is now validated}
\begin{itemize}
  \item Replay-side reproducibility is stable (schema checks, lineage, fixture snapshots).
  \item Real-call execution is operational after key-path recovery (HTTP auth failures removed in full-range runs).
  \item Failure concentration is now clearly behavioral: keyword reversion and normalization-parse breakdown.
\end{itemize}

\subsection{What remains limited}
\begin{itemize}
  \item All four tested model settings currently show 0 pass under strict judge criteria.
  \item Cross-condition prompt design is still narrow (few-shot/repair-loop variants not yet fully explored).
\end{itemize}

\subsection{Methodological implication}
For constrained-language evaluation, provenance and schema controls are necessary but not sufficient. Once infra is stable, model behavior can still collapse under strict rule adherence. Therefore, next experiments should explicitly separate prompt-design effects, decoding constraints, and judge strictness sensitivity.

\section{Conclusion}
We presented a reproducible framework to study how coding LLMs behave under Python-like confusion grammar. Initial evidence suggests substantial reversion pressure toward canonical Python. The framework is ready for larger-scale experiments, and future iterations will emphasize stronger statistical evidence and educational impact analysis.


\section{Reproducibility and artifact availability}
All scripts, logs, and metrics are versioned in the repository. Representative paths:
\begin{itemize}
  \item Tooling: \texttt{ocaml-confusion-lang/}
  \item Results: \texttt{docs/research/results/}
  \item Logs: \texttt{docs/research/log/}
  \item Status snapshot: \texttt{docs/research/context-state.json}
  \item Paper draft: \texttt{docs/research/papers/v2.md}
\end{itemize}

\bibliographystyle{plainnat}
\bibliography{refs}

\end{document}
